%! Characteristics of Exponential Functions — Unit 6, Lesson 6.0 — Group Activity (Answer Key)
\documentclass[10pt]{article}
\usepackage{algebra2-article}
\usepackage{algebra2-key}   % pulls in -boxes; do NOT also load -boxes
\newcommand{\TallMath}[1]{$\displaystyle #1\rule[-1.4em]{0pt}{3.2em}$}
\newcommand{\lgrid}[4]{%
  \draw[step=1, linegray!40, very thin] (#1,#3) grid (#2,#4);
  \draw[->, semithick, charcoal] (#1,0)--(#2,0) node[below right, font=\scriptsize]{$x$};
  \draw[->, semithick, charcoal] (0,#3)--(0,#4) node[left, font=\scriptsize]{$y$};}
% #2 is ln(b):  ln2 = 0.693147   ln3 = 1.098612   ln4 = 1.386294
\newcommand{\expcurve}[4]{\draw[very thick, forest, domain=#3:#4, samples=120, smooth]
  plot (\x,{#1*exp((#2)*(\x))});}

\begin{document}

\pageheader{Unit 6, Lesson 6.0}{Group Activity --- Answer Key}
\namepartnerperiod

\vspace{0.08in}

\begin{headlinebox}{forestbg}
\small\textbf{Reading Exponential Graphs.} An exponential graph never turns around, never breaks, and
never crosses the $x$-axis --- so the questions worth asking are \emph{which way does it run} and
\emph{which line is it sneaking up on}. With your partner, read each graph carefully and
\emph{justify} your answers. Write every domain and range in \textbf{interval notation}, and check
each one against the equation.
\end{headlinebox}

\vspace{0.06in}

\begin{tcolorbox}[colback=white, colframe=black!40, title=\textbf{Tier R --- Remediate},
    fonttitle=\bfseries, arc=1.5mm, left=3mm, right=3mm, top=2mm, bottom=2mm, breakable]
The graph shows $g(x) = 3^{x}$.
\begin{center}
\begin{tikzpicture}[scale=0.44]
  \lgrid{-3}{3}{-1}{10}
  \begin{scope}
    \clip (-3,-1) rectangle (3,10);
    \expcurve{1}{1.098612}{-3}{2.09}
  \end{scope}
  \draw[navy, dashed, semithick] (-3,0)--(3,0);
  \node[navy, font=\scriptsize, above left] at (-1.0,0) {$y=0$};
  \fill[charcoal] (0,1) circle (3pt) node[left, font=\scriptsize]{$(0,1)$};
  \fill[charcoal] (1,3) circle (3pt) node[right, font=\scriptsize]{$(1,3)$};
  \fill[charcoal] (2,9) circle (3pt) node[left, font=\scriptsize]{$(2,9)$};
\end{tikzpicture}
\end{center}
\begin{enumerate}[leftmargin=1.7em, itemsep=7pt]
  \item Write $g$ in the form $ab^{x}$: \; $a=$ \ans{1} and $b=$ \ans{3}. \\
        Since $b$ is \ans{greater} than $1$, this is exponential \ans{growth}.
  \item \textbf{Domain:} \ans{$(-\infty,\infty)$} \quad \textbf{Range:} \ans{$(0,\infty)$}
  \item \textbf{$y$-intercept:} \ans{$(0,1)$}. \; Check it against $a$: are they the same?
        \ans{yes}
  \item Try to find an \textbf{$x$-intercept} by asking when $3^{x}=0$. From the table you built in
        the Warm-Up, $3^{-1}=$ \ans{$\frac13$} and $3^{-4}=$ \ans{$\frac{1}{81}$}. \\
        Do these ever reach $0$? \ans{no} \; So this graph has \ans{no}
        $x$-intercept.
  \item $g$ is \textbf{increasing} on \ans{$(-\infty,\infty)$}. \quad Its absolute extrema:
        \ans{none}
  \item \textbf{Horizontal asymptote:} \ans{$y=0$}. \; The graph approaches it on the
        \ans{left} side.
  \item \textbf{End behavior:} as $x\to+\infty$, $g(x)\to$ \ans{$+\infty$}. \\
        As $x\to-\infty$, $g(x)\to$ \ans{$0$} --- careful, this is \emph{not} $-\infty$.
        Explain in one sentence what the graph is actually doing on the left. \ans{It flattens out
        and hugs the line $y=0$ from above, getting closer forever without ever touching it.}
\end{enumerate}
\end{tcolorbox}

\begin{tcolorbox}[colback=white, colframe=black!40, title=\textbf{Tier A --- Approaching Proficiency},
    fonttitle=\bfseries, arc=1.5mm, left=3mm, right=3mm, top=2mm, bottom=2mm, breakable]
Two exponential functions. Fill in the whole table, then answer the two justification questions ---
and pay attention to which rows come out \emph{identical}.
\begin{center}
\begin{tikzpicture}[scale=0.40]
  % p(x) = 4^x
  \begin{scope}
    \lgrid{-3}{3}{-1}{10}
    \begin{scope}
      \clip (-3,-1) rectangle (3,10);
      \expcurve{1}{1.386294}{-3}{1.66}
    \end{scope}
    \draw[navy, dashed, semithick] (-3,0)--(3,0);
    \fill[charcoal] (0,1) circle (3pt);
    \fill[charcoal] (1,4) circle (3pt);
    \node[charcoal, font=\scriptsize] at (0,-2.6) {$p(x)=4^{x}$};
  \end{scope}
  % q(x) = 8(1/4)^x
  \begin{scope}[xshift=10cm]
    \lgrid{-1}{5}{-1}{10}
    \begin{scope}
      \clip (-1,-1) rectangle (5,10);
      \expcurve{8}{-1.386294}{-0.16}{5}
    \end{scope}
    \draw[navy, dashed, semithick] (-1,0)--(5,0);
    \fill[charcoal] (0,8) circle (3pt);
    \fill[charcoal] (1,2) circle (3pt);
    \node[charcoal, font=\scriptsize] at (2,-2.6) {$q(x)=8\left(\tfrac14\right)^{x}$};
  \end{scope}
\end{tikzpicture}
\end{center}
\vspace{0.14in}
\renewcommand{\arraystretch}{1.5}
\begin{tabularx}{\linewidth}{@{}>{\bfseries}p{0.30\linewidth} X X@{}}
 & \textbf{$p(x)=4^{x}$} & \textbf{$q(x)=8\left(\frac14\right)^{x}$} \\
\hline
Initial value $a$ & \ans{1} & \ans{8} \\
Base $b$ & \ans{4} & \ans{$\frac14$} \\
Growth or decay? & \ans{growth} & \ans{decay} \\
Domain & \ans{$(-\infty,\infty)$} & \ans{$(-\infty,\infty)$} \\
Range & \ans{$(0,\infty)$} & \ans{$(0,\infty)$} \\
$y$-intercept & \ans{$(0,1)$} & \ans{$(0,8)$} \\
$x$-intercept & \ans{none} & \ans{none} \\
Increasing or decreasing? & \ans{increasing} & \ans{decreasing} \\
Horizontal asymptote & \ans{$y=0$} & \ans{$y=0$} \\
Asymptote approached on the\dots & \ans{left} & \ans{right} \\
Absolute max or min? & \ans{neither} & \ans{neither} \\
\end{tabularx}
\vspace{4pt}

\textbf{Justify 1.} One of these functions rises and the other falls. Point to the exact part of
each \emph{equation} that decides it, and explain why that part has the effect it does. (Be careful:
the initial values differ too --- is $a$ what decides it?) \ansline{The \textbf{base} decides it. For
$p$, $b=4>1$: every time $x$ goes up by $1$ the output is multiplied by $4$, so it grows. For $q$,
$b=\frac14$, which is between $0$ and $1$: every step multiplies by a proper fraction, so the output
shrinks. $a$ is \emph{not} what decides it --- $a$ only sets the starting height ($y$-intercept), and
both of these have $a>0$. If $a$ controlled direction, $q$ would rise, since its $a$ is the larger
one.}

\textbf{Justify 2.} Five rows of your table came out \emph{identical} for both functions. List them,
and explain what those five rows tell you about \emph{every} function of the form $y=ab^{x}$ with
$a>0$. \ansline{Domain, range, $x$-intercept, horizontal asymptote, and absolute max/min are the
same for both. So for \emph{any} $y=ab^{x}$ with $a>0$: the domain is all reals (every real number is
a legal exponent), the range is $(0,\infty)$, there is never an $x$-intercept, the horizontal
asymptote is always $y=0$, and there are never any absolute extrema. Only three things change from
one such function to the next --- the starting height $a$, whether it rises or falls, and which side
of the graph hugs the asymptote.}
\end{tcolorbox}

\begin{tcolorbox}[colback=white, colframe=black!40, title=\textbf{Tier E --- Extension},
    fonttitle=\bfseries, arc=1.5mm, left=3mm, right=3mm, top=2mm, bottom=2mm, breakable]
\textbf{Part 1 --- No graph allowed.} Fill in each row from the equation alone.
\begin{center}
\renewcommand{\arraystretch}{1.6}
\begin{tabularx}{\linewidth}{@{}l X X X X@{}}
\hline
\textbf{Function} & \textbf{Base $b$} & \textbf{Is $b>1$ or $0<b<1$?} & \textbf{$y$-intercept} & \textbf{Range} \\
\hline
$y=7(1.5)^{x}$ & \ans{$1.5$} & \ans{$b>1$} & \ans{$(0,7)$} & \ans{$(0,\infty)$} \\
$y=7(0.5)^{x}$ & \ans{$0.5$} & \ans{$0<b<1$} & \ans{$(0,7)$} & \ans{$(0,\infty)$} \\
$y=-2\cdot 3^{x}$ & \ans{$3$} & \ans{$b>1$} & \ans{$(0,-2)$} & \ans{$(-\infty,0)$} \\
$y=\left(\frac54\right)^{x}$ & \ans{$\frac54$} & \ans{$b>1$} & \ans{$(0,1)$} & \ans{$(0,\infty)$} \\
\hline
\end{tabularx}
\end{center}
Exactly one of these four has a range that is \emph{not} $(0,\infty)$. Which one, and what in the
equation tells you so \emph{before} you graph anything? \ansline{$y=-2\cdot 3^{x}$. The base is
still positive, so $3^{x}$ is always positive --- but $a=-2$ is \emph{negative}, and a negative times
a positive is negative for every $x$. So every output is below the axis and the range is
$(-\infty,0)$. The horizontal asymptote is still $y=0$; the graph just approaches it from
\emph{underneath}. The sign of $a$ decides which side of the asymptote the whole graph lives on.}

\vspace{5pt}
\textbf{Part 2 --- A model, and the line it never crosses.} A patient takes a $200$-milligram dose of
a medication. Each hour, the body clears some of it and about $70\%$ of the previous hour's amount
remains, so the amount in the bloodstream after $t$ hours is
\[
  A(t) = 200(0.7)^{t} \ \text{ milligrams.}
\]
\begin{center}
\renewcommand{\arraystretch}{1.25}
\begin{tabular}{|c|c|c|c|c|c|}
\hline
$t$ (hours) & $0$ & $1$ & $2$ & $3$ & $4$ \\ \hline
$A(t)$ (mg) & $200$ & $140$ & $98$ & $68.6$ & $48.0$ \\ \hline
\end{tabular}
\end{center}
\begin{enumerate}[label=(\alph*), leftmargin=1.9em, itemsep=6pt]
  \item Identify $a=$ \ans{200} and $b=$ \ans{$0.7$}. Is this growth or decay?
        \ans{decay} \\
        What does each of those two numbers \emph{mean} for this patient? \ansline{$a=200$ is the
        dose at the moment it is taken --- the amount at $t=0$. $b=0.7$ means $70\%$ of whatever is
        present survives each hour, so $30\%$ is cleared each hour.}
  \item \textbf{Confirm the fingerprint.} $\dfrac{140}{200}=$ \ans{$0.7$} and
        $\dfrac{98}{140}=$ \ans{$0.7$}. \; What is this constant called, and why does finding it
        prove the model is exponential rather than linear? \ansline{It is the \textbf{constant
        ratio}, and it equals $b$. A linear model would have a constant \emph{difference}; here the
        differences are $60$, $42$, $29.4$ --- not constant --- while every ratio is $0.7$.}
  \item Find $A(3)$ from the table and write one sentence explaining what it means in context.
        \ansline{$A(3)=68.6$: three hours after the dose is taken, about $68.6$ mg of the
        medication remain in the patient's bloodstream.}
  \item What is the domain of $A$ \emph{algebraically} --- that is, which exponents are legal?
        \ans{all real numbers} \\
        Does the context agree, or does it restrict things further? Explain. \ansline{The context
        restricts it. Negative $t$ would describe time before the dose was taken, when the model
        does not apply, so the sensible domain is $t\ge0$, or $[0,\infty)$.}
  \item The graph of $A$ has horizontal asymptote $y=0$, and the model has \textbf{no}
        $t$-intercept. State what that means about the drug, then say whether you believe it.
        \ansline{The model says the amount keeps shrinking toward $0$ but is never exactly $0$ ---
        the drug is never \emph{completely} gone, no matter how long you wait.\\[3pt]
        That is not literally true. A real dose is a finite number of molecules, so eventually the
        last one is cleared and the amount really does reach $0$. The model is accurate for the
        first several hours and stops being meaningful in the far tail, once the predicted amount is
        smaller than a single molecule.}
  \item A classmate says ``the patient loses $60$ mg every hour --- look, $200$ to $140$.'' Use the
        table to show that this is wrong, and name what is actually constant instead.
        \ansline{The next drops are $140-98=42$ and $98-68.6=29.4$, not $60$ --- the amount lost
        shrinks every hour. What stays constant is the \emph{ratio}: the patient loses $30\%$ of
        whatever is still there, and $30\%$ of a smaller amount is a smaller loss.}
\end{enumerate}
\end{tcolorbox}

\begin{teachernote}
\textbf{Tier R, item 7} is the item to check first when you circulate. ``$g(x)\to-\infty$'' on the
left is the reflex answer from Unit 3's odd-degree polynomials, and it is wrong for a reason worth
saying out loud: an exponential graph's left arm does not run \emph{down}, it runs \emph{flat}. Pair
it with the same student's item 6 --- if they wrote the asymptote correctly and then contradicted it
one line later, the fix is one sentence, not a re-teach.

\textbf{Tier A, Justify 2} is the real assessment of the day (\textbf{A2.F.2b}). Students who list
the five shared rows have understood that domain, range, asymptote, ``no zeros,'' and ``no extrema''
are \emph{family} properties, not facts about one graph. Groups that instead list rows that differ
have read the table but not the point; send them back with ``which of these would still be true if I
changed the base to $10$?''

\textbf{Tier E, Part 1, row 3} is deliberately the only negative-$a$ function students meet today.
Do \emph{not} let it become a lesson on reflections --- that is 6.2. The takeaway is narrower and
worth exactly one sentence: the sign of $a$ says which side of the asymptote the graph lives on.
Watch for ``$b=-2$'' here; the negative belongs to $a$, and the base is still $3$.

\textbf{Tier E, Part 2(e)} rewards honesty about models, not a formula. Full credit needs both
halves: what the mathematics claims, \emph{and} the observation that a finite pile of molecules must
eventually run out. Part 2(f) is the constant-ratio idea in its most useful disguise and is worth
sharing out to the whole class --- ``$30\%$ of what is left'' is how every decay context in Lessons
6.1 and 6.4 will be worded.
\end{teachernote}

\end{document}
